\chapter{Códigos de términos y fórmulas}
\label{cap:codigos}

Con listas dentro de la aritmética, ya se puede codificar sintaxis. Este capítulo construye
\(\codigo{\varphi}\): el término del lenguaje objeto que denota la fórmula \(\varphi\).

\section{Dos codificaciones, y por qué hay dos}

El proyecto tiene una codificación \emph{de juguete} y una \emph{real}, y conviene no confundirlas.

La primera —el nivel B— codifica cadenas de símbolos de un alfabeto finito de doce elementos:

\leanfrag{Sym}
\leanfrag{gNat}
\leanfrag{encode}

\noindent Es la construcción clásica de Gödel, útil para explicar la idea y para el teorema de
inyectividad, pero no sirve para el argumento real: codifica \emph{cadenas}, no el árbol sintáctico
de las fórmulas del kernel.

\section{La codificación real: por constructor}

La segunda codificación —el nivel C— recorre la estructura de \ident{Term} y \ident{Formula} y
asigna una etiqueta a cada constructor.

\leanfrag{termCode}
\leanfrag{formCode}

\begin{matematica}[Las etiquetas]
\[
\begin{array}{llll}
\bot \mapsto 2 & \text{átomo} \mapsto 3 & {\doteq} \mapsto 4 & {\Rightarrow} \mapsto 5\\
\forall \mapsto 6 & \wedge \mapsto 7 & \vee \mapsto 8 & \exists \mapsto 9
\end{array}
\]
y, en los términos, \(0\) para variables y \(1\) para aplicaciones. Un código es una lista cuya
cabeza es la etiqueta y cuya cola son los códigos de las partes.
\end{matematica}

\begin{leccion}[Dónde vive cada cosa]
\ident{formCode} es una función \textbf{de Lean}: toma un valor del tipo \ident{Formula} y devuelve
un valor del tipo \ident{Term}. Su \emph{resultado} es un término del lenguaje objeto. Por eso la
chapa dice \emph{Lean} sobre \emph{código}, y por eso el capítulo de nomenclatura insiste: la
función que calcula el código es meta; el código es objeto.
\end{leccion}

\section{Inyectividad, sin suponer consistencia}

Que la codificación sea inyectiva es lo que la hace utilizable: dos fórmulas distintas tienen
códigos distintos, luego un código determina su fórmula.

\begin{teorema}[La codificación es inyectiva]
Si \(\codigo{\varphi} = \codigo{\psi}\) entonces \(\varphi = \psi\).
\end{teorema}

\leanfrag{formCode_injective}

\begin{leccion}[Un detalle que importa más de lo que parece]
Este teorema se demuestra \textbf{en la metateoría}, por inducción estructural, y \textbf{no supone
la consistencia} de la teoría objeto. Es un hecho sobre los constructores de un tipo inductivo de
Lean, no sobre lo que la aritmética demuestre. Si dependiera de la consistencia, todo el edificio
sería circular.
\end{leccion}

\section{De códigos a demostrabilidad}

Con la codificación se pueden definir, ya en la metateoría, los dos predicados que dan sentido a
todo lo demás:

\leanfrag{IsFormula}
\leanfrag{Provable}

\noindent Y el primer resultado que los conecta:

\leanfrag{provable_formCode_iff}

\begin{muro}[Esto todavía no es Gödel]
\ident{Provable} es un predicado \emph{de Lean} sobre términos: dice que existe una fórmula con ese
código y que la teoría la demuestra. Es correcto y es inútil para la incompletitud, porque la
teoría objeto no puede hablar de él — vive en la metateoría. Lo que hace falta es una \emph{fórmula
del lenguaje objeto} que exprese «existe una demostración». Ése es el capítulo siguiente.
\end{muro}
